%! AP Statistics — Unit 1, Lesson 1.6 — Guided Notes
\documentclass[10pt]{article}
\usepackage{apstats-article}
\usepackage{apstats-boxes}

%=============================================================================
\begin{document}

\pageheader{Unit 1, Lesson 1.6}{Guided Notes}
\namedateperiod

\vspace{0.12in}

% ── OBJECTIVE ────────────────────────────────────────────────────────────────
\begin{objectivebox}
\small By the end of this lesson, I will be able to\dots\\[4pt]
\writeline\\[1pt]
\writeline\\[1pt]
\writeline
\end{objectivebox}

\vspace{0.1in}

% ── VOCABULARY ───────────────────────────────────────────────────────────────
\begin{vocabbox}
\termblanklong{SOCS}
\termblanklong{Shape}
\termblanklong{Skewed right (positive skew)}
\termblanklong{Skewed left (negative skew)}
\termblanklong{Symmetric}
\termblanklong{Unimodal / Bimodal / Uniform}
\termblanklong{Outlier}
\end{vocabbox}

\vspace{0.1in}

% ── HOOK ─────────────────────────────────────────────────────────────────────
\begin{hookbox}
\small\textbf{Two Histograms --- Describe What You See}

\vspace{8pt}
\textbf{Histogram A:} U.S.\ household incomes. Most households earn below
\$80{,}000; a small number earn \$200{,}000 or more; the distribution has a
very long tail toward higher values.

\vspace{6pt}
\textbf{Histogram B:} Heights of adult women. The bars form a roughly
symmetric, bell-shaped mound centered near 64 inches; no extreme outliers.

\vspace{8pt}
\begin{enumerate}[leftmargin=1.5em, itemsep=6pt]
  \item Describe Histogram A in one word or phrase.\\[5pt]
        \writeline

  \item Describe Histogram B in one word or phrase.\\[5pt]
        \writeline

  \item Why might the shape of a distribution matter for how we
        summarize it?\\[5pt]
        \writeline\\[1pt]\writeline
\end{enumerate}
\end{hookbox}

\vspace{0.1in}

% ── SOCS FRAMEWORK ───────────────────────────────────────────────────────────
\begin{notesbox}{The SOCS Framework --- Describing Any Quantitative Distribution}
\small

\begin{multicols}{2}

\textbf{\textcolor{navy}{S --- Shape}}
\begin{itemize}[itemsep=2pt]
  \item Overall pattern: \blank{2.5cm}, skewed right, skewed left
  \item Number of peaks: \blank{1.5cm} (one peak) or \blank{1.5cm} (two peaks)
  \item Special case: \blank{1.8cm} (all bars roughly equal height)
\end{itemize}

\vspace{4pt}
\textbf{\textcolor{navy}{O --- Outliers}}
\begin{itemize}[itemsep=2pt]
  \item Values that fall \blank{3cm} from the overall pattern
  \item Note their approximate \blank{2cm} and how far they are from the bulk
  \item Always \blank{2cm} outliers; do not automatically remove them
\end{itemize}

\columnbreak

\textbf{\textcolor{navy}{C --- Center}}
\begin{itemize}[itemsep=2pt]
  \item Informal estimate of a \blank{2.5cm} value
  \item For now: use the \blank{2cm} or the approximate balance point
  \item Lesson 1.7 will give us the \blank{1.5cm} and \blank{1.5cm}
\end{itemize}

\vspace{4pt}
\textbf{\textcolor{navy}{S --- Spread}}
\begin{itemize}[itemsep=2pt]
  \item How \blank{2cm} or \blank{2cm} the distribution is
  \item For now: describe the \blank{1.8cm} (minimum to maximum)
  \item Lesson 1.7: IQR and standard deviation
\end{itemize}

\end{multicols}

\vspace{4pt}
\textbf{\textcolor{navy}{AP Context Rule:}} Every sentence must name the
\blank{2.5cm} and the \blank{2cm}.
Example: ``The distribution of \underline{\hspace{3cm}} for these
\underline{\hspace{2.5cm}} is\ldots''

\end{notesbox}

\vspace{0.1in}

% ── SHAPE GUIDE ──────────────────────────────────────────────────────────────
\begin{notesbox}{Shape Vocabulary --- Recognizing Common Patterns}
\small

\renewcommand{\arraystretch}{2.0}
\begin{tabularx}{\linewidth}{@{} >{\bfseries\color{navy}}p{0.22\linewidth}
                                   p{0.38\linewidth}
                                   X @{}}
\rowcolor{sky}
\textbf{Shape} & \textbf{What it looks like} & \textbf{Real-world example} \\
\midrule
Symmetric      & Mirror image on each side; peak in the center & \blank{3.5cm} \\
Skewed right   & Long tail toward \blank{1.5cm} values; peak on the left & \blank{3.5cm} \\
Skewed left    & Long tail toward \blank{1.5cm} values; peak on the right & \blank{3.5cm} \\
Unimodal       & \blank{3.5cm} clear peak & \blank{3.5cm} \\
Bimodal        & \blank{2cm} clear peaks & \blank{3.5cm} \\
Uniform        & All bars roughly \blank{2cm} height & \blank{3.5cm} \\
\end{tabularx}

\vspace{6pt}
\textbf{Common error:} ``Skewed right'' does \textit{not} mean the peak is on
the right. It means the \blank{2cm} is on the right. Remember: the skew
direction = the direction the \blank{1.5cm} points.

\end{notesbox}

\newpage

% ── WORKED EXAMPLE ───────────────────────────────────────────────────────────
\begin{notesbox}{Worked Example --- SAT Math Scores for 200 Students}
\small

\noindent The histogram shows SAT Math scores (out of 800) for 200 students.
Key features: roughly symmetric bell shape; center near 520; scores range from
about 300 to 750; two students scored above 750 (possible outliers).

\vspace{8pt}
\textbf{Model SOCS response (label each sentence with S, O, C, or S):}

\vspace{6pt}
\begin{tcolorbox}[colback=white, colframe=slate, boxrule=0.5pt, arc=2mm,
    left=3mm, right=3mm, top=3mm, bottom=3mm]
\textit{``The distribution of SAT Math scores for these 200 students is
approximately symmetric and unimodal, with no obvious skew.
[~\blank{0.5cm}~] There appear to be two potential outliers above 750,
which stand out from the rest of the distribution. [~\blank{0.5cm}~]
The center of the distribution is approximately 520 points, which represents
a typical score for these students. [~\blank{0.5cm}~] The spread is roughly
300 to 750 points, giving a range of about 450 points.''}
[~\blank{0.5cm}~]
\end{tcolorbox}

\vspace{8pt}
\textbf{What made this a strong AP response?}
\begin{enumerate}[leftmargin=1.5em, itemsep=3pt]
  \item Context (\textit{SAT Math scores for these 200 students}) appears in
        the \blank{2.5cm} sentence and is referenced throughout.
  \item All \blank{1cm} SOCS components are present.
  \item Outliers were \blank{2cm} (``appear to be'') rather than stated as
        definite --- we will learn the formal rule in Lesson 1.7.
  \item Center and spread include \blank{2cm} with appropriate labels.
\end{enumerate}

\end{notesbox}

\vspace{0.1in}

% ── GUIDED PRACTICE ──────────────────────────────────────────────────────────
\begin{practicebox}
\small
\textbf{A survey records the number of text messages sent per day by
40 high school students.}

Key features from the dotplot: most students send 20--60 messages; a small
cluster sends 100--140 messages; one student sends 180 messages (far from the
rest); center near 45 messages; values range from 5 to 180.

\vspace{8pt}
\begin{enumerate}[leftmargin=1.5em, itemsep=8pt]
  \item What is the shape of the distribution? Is it unimodal or bimodal?
        Explain what features in the description support your answer.\\[6pt]
        \writeline\\[1pt]\writeline

  \item Identify any potential outliers. What values qualify and why?\\[6pt]
        \writeline\\[1pt]\writeline

  \item Estimate the center and describe the spread. Use appropriate units.\\[6pt]
        \writeline\\[1pt]\writeline

  \item Write a complete SOCS description of the distribution of daily text
        messages for these 40 students. Include context in every sentence.\\[6pt]
        \writeline\\[1pt]\writeline\\[1pt]\writeline\\[1pt]\writeline
\end{enumerate}
\end{practicebox}

\vspace{0.1in}

% ── SPIRAL / BIG IDEAS ────────────────────────────────────────────────────────
\begin{spiralbox}
\small
\begin{multicols}{2}

\textbf{This lesson connects forward to\dots}
\begin{itemize}[itemsep=2pt]
  \item \textbf{Lesson 1.7:} Mean and median formalize center; range, IQR,
        and standard deviation formalize spread.
  \item \textbf{Lesson 1.9:} Compare \textit{two} distributions using SOCS ---
        same framework, two groups, comparative language.
  \item \textbf{AP Exam:} Free-response almost always asks you to describe a
        distribution --- SOCS earns full credit when done completely with context.
\end{itemize}

\columnbreak

\textbf{Big Idea --- Write in your own words:}\\
\textit{Why is it not enough to say ``the distribution is skewed right''?
What else does the AP grader need to see?}\\[2pt]
\writeline\\[1pt]\writeline\\[1pt]\writeline

\end{multicols}
\end{spiralbox}

\vspace{0.1in}

\begin{notesbox}{Additional Notes}
\vspace{0pt}
\writeline\\\writeline\\\writeline\\\writeline\\\writeline\\\writeline
\end{notesbox}

\end{document}
